% figures/chart-validation.tex  —  \input in Ch.4 (faithfulness check)
% Our pairs-F1 vs Chen 2016 published pairs-F1; points on y=x => protocol-faithful.
% Numbers from Table~\ref{tab:flickr-faith}.
% Requires: pgfplots + sv* palette (sv base style).

\begin{figure}[ht]
\centering
\begin{tikzpicture}
\begin{axis}[
  sv base,
  width=0.66\linewidth, height=0.66\linewidth,
  xmin=0.20, xmax=0.55, ymin=0.20, ymax=0.55,
  xlabel={Chen 2016 published pairs-F1},
  ylabel={our pairs-F1},
  grid=both, minor tick num=1,
  axis lines=left,
  legend pos=north west,
  legend cell align=left,
]
% y = x reference
\addplot[domain=0.20:0.55, samples=2, dashed, svGray, line width=0.9pt]{x};
\addlegendentry{$y=x$ (perfect match)}
% Random
\addplot[only marks, mark=*, mark size=2.6pt, color=svGray]
  coordinates {(0.310,0.298)(0.304,0.301)(0.320,0.301)(0.261,0.270)(0.248,0.227)};
\addlegendentry{Random}
% PoiPopularity
\addplot[only marks, mark=square*, mark size=2.6pt, color=svTeal]
  coordinates {(0.384,0.443)(0.365,0.413)(0.507,0.510)(0.436,0.439)(0.316,0.320)};
\addlegendentry{PoiPopularity}
\end{axis}
\end{tikzpicture}
\caption{Protocol-faithfulness check: our reproductions of the assumption-free baselines fall
on the $y=x$ line of Chen 2016 (each marker is one Flickr city; cf.\
Table~\ref{tab:flickr-faith}). \texttt{Random} in particular has no modelling choices, so its
agreement confirms that our data, protocol, and metric match the literature.}
\label{fig:chart-validation}
\end{figure}
